\scanpage{47}
% L1
d'après le n\textsuperscript{o} précédent que \deleted{la composante \unsure{$j$-bipure|$j$-bipures} de} \inserted{l'image dans}
\[
\underset{\substack{U\,\text{ouvert}\\\deleted{\text{\unsure{de}}}\\\text{dense}}}{\underrightarrow{\lim}}H^j(U,\mathbb Q(0))
\]
de $H^j(X,\mathbb Z(0))$.

% L4
En fait, on constate (via dualité) que si $X$ est lisse (sans doute inutile) de dim $\leq n$, et si $i\leq n$ alors pour tt ouvert dense $U$ de $X$,
\[
H^i(X)\mathrel{\unsure{-|\to}}H^i(U)
\]
est \uline{\deleted{bijectif}\inserted{injectif} mod} $D_{i-2}(M^+(X))$ et \uline{bijectif mod} $\underline{D_{i-1}(M^+(X))}$, ce qui donne un invariant birationnel dans
\[
D_{i}(M^+(X))/D^{\unsure{+}}_{i-1}(M^+(X)),
\]
lequel est d'ailleurs dans l'image de $M^+_{(i)}(X)$, donc définit un objet de
\[
M^+_{(i)}(X)/[M^+_{(i-1)}(X)\unsure{;|)|,}\deleted{\unsure{+}}\xi M^+_{(i-2)}(X)],
\]
donc un élément de $\mathbb Z^{\sigma_i(X)}$. C'est l'invariant birationnel $\gamma_i(X)$.
% archive sheet number: 46
